\chapter{Tinjauan Pustaka}
\label{chap:tinjauan-pustaka}

% Ubah bagian-bagian berikut dengan isi dari tinjauan pustaka

Demi mendukung penelitian ini, bab ini akan memaparkan hasil penelitian terdahulu yang relevan serta landasan teori yang melatarbelakangi metode yang diusulkan.

\section{Hasil Penelitian Terdahulu}
\label{sec:penelitian-terdahulu}

Beberapa penelitian terdahulu yang relevan dengan topik deteksi konjungtivitis menggunakan metode \textit{machine learning} dan \textit{deep learning} dirangkum dalam Tabel \ref{tab:penelitian-terdahulu}.

\begin{longtable}{|c|p{4.5cm}|p{8.5cm}|}
  \caption{Penelitian Terdahulu}\label{tab:penelitian-terdahulu} \\
  \hline
  \textbf{No} & \textbf{Penulis \& Tahun} & \textbf{Pembahasan} \\
  \hline
  \endfirsthead
  
  \multicolumn{3}{c}{\tablename\ \thetable\ -- \textit{Lanjutan dari halaman sebelumnya}} \\
  \hline
  \textbf{No} & \textbf{Penulis \& Tahun} & \textbf{Pembahasan} \\
  \hline
  \endhead
  
  \hline
  \endfoot
  
  \hline
  \endlastfoot
  
  1 & \multicolumn{2}{p{13.3cm}|}{\textit{Automated Conjunctivitis Detection Using Xception Model with Transfer Learning}} \\
  \cline{2-3}
    & Syed Sofiya Ali dan Suman Kumar Swarnkar, 2025 \parencite{ali2025automated} & 
    Penelitian ini mengklasifikasi konjungtivitis biner menggunakan model Xception CNN standar dan transfer learning yang dilatih pada 582 citra mata. Model dilatih dalam dua tahap (pembekuan lapisan dan fine-tuning) dengan augmentasi citra RGB standar. Model mencapai akurasi pengujian sebesar 95,16\%, AUC-ROC 0,9484, precision 0,9773, recall 0,9412, dan F1-Score 0,9143. Kelemahannya adalah klasifikasi dilakukan secara global pada seluruh citra mata tanpa segmentasi sklera, sehingga rentan terganggu oleh derau non-okular (kelopak/bulu mata). \\
  \hline

  2 & \multicolumn{2}{p{13.3cm}|}{\textit{Enhanced Conjunctivitis Detection: Leveraging ResNet and Spatial Attention Mechanisms for Superior Accuracy}} \\
  \cline{2-3}
    & S. N. Sangeethaa dkk., 2024 \parencite{sangeethaa2024enhanced} & 
    Penelitian ini mengusulkan metode deteksi konjungtivitis berbasis arsitektur ResNet dengan integrasi Spatial Attention Mechanism. Diuji pada dataset berisi 358 citra mata (181 sehat dan 177 sakit), model memperoleh akurasi sebesar 92,5\%, sensitivitas 90,7\%, spesifisitas 94,3\%, dan F1-Score 0,92. Kelemahannya adalah belum mengintegrasikan prapemrosesan segmentasi atau lokalisasi untuk membuang derau non-okular di sekitar area mata. \\
  \hline

  3 & \multicolumn{2}{p{13.3cm}|}{\textit{Automatic diagnosis of keratitis using object localization combined with cost-sensitive deep attention convolutional neural network}} \\
  \cline{2-3}
    & Jiewei Jiang dkk., 2023 \parencite{jiang2023automatic} & 
    Penelitian ini mengusulkan kombinasi lokalisasi objek Single Shot Multibox Detector (SSD) untuk memotong area kornea-konjungtiva (Conj\_Cor) dan pengklasifikasi Cost-Sensitive Deep Attention CNN (CDACNN). Diuji pada 12.407 citra slit-lamp, model ini mencapai AUC hingga 0,998. Kelemahannya adalah model difokuskan pada keratitis kornea dan membutuhkan anotasi koordinat kotak pembatas yang rumit. Namun, lokalisasi objek ini sangat efektif menghilangkan derau kelopak mata. \\
  \hline

  4 & \multicolumn{2}{p{13.3cm}|}{\textit{A Deep Learning-Based Segmentation Technique for Automatic Detection of Conjunctivitis from Conjunctival Eye Images}} \\
  \cline{2-3}
    & S. Madhusudhan dan S. Anitha, 2026 \parencite{madhusudhan2026deep} & 
    Penelitian ini menggunakan segmentasi Attention-U-Net++ dipadukan dengan Swin/ViT Transformer untuk deteksi otomatis konjungtivitis berdasarkan fitur densitas pembuluh darah (vessel density index) dan indeks kemerahan sklera. Model mencapai akurasi segmentasi 90,3\%, spesifisitas 0,93, dan F1-Score 0,89 pada dataset citra mata. Kelemahannya adalah kompleksitas model komputasi yang tinggi karena penggabungan CNN dan Transformer. \\
  \hline

  5 & \multicolumn{2}{p{13.3cm}|}{\textit{Illumination-Robust Conjunctival Image Preprocessing for Accurate Segmentation and Anemia Detection Using Deep Learning}} \\
  \cline{2-3}
    & Jose Humberto Fuentes-Beingolea dkk., 2025 \parencite{fuentes2025illumination} & 
    Penelitian ini merancang pipeline prapemrosesan citra kebal cahaya (normalisasi histogram, inpainting glare berbasis ruang warna LAB, dan CLAHE) untuk deteksi anemia non-invasif pada konjungtiva mata. Dengan preprocessing tersebut, presisi segmentasi U-Net naik dari 80,08\% menjadi 84,22\%, dan F1-Score klasifikasi CNN-ResNet naik dari 81,91\% menjadi 89,15\% pada kondisi pencahayaan liar. Kelemahannya adalah inpainting glare yang terlalu luas dapat mengaburkan detail pembuluh darah asli di bawahnya. \\
  \hline

\end{longtable}

Berdasarkan hasil pembandingan di atas, posisi penelitian yang diusulkan adalah menggabungkan keunggulan dari segmentasi sklera otomatis untuk menghilangkan \textit{noise} non-okular (terinspirasi dari \textcite{jiang2023automatic}), perbaikan cahaya dan restorasi \textit{glare} (terinspirasi dari \textcite{fuentes2025illumination}), penguatan fitur kemerahan pembuluh darah (\textit{vessel enhancement} menggunakan CLAHE terinspirasi oleh \textcite{madhusudhan2026deep}), serta pemanfaatan arsitektur \textit{Attention-driven CNN} berbasis ResNet \parencite{sangeethaa2024enhanced} untuk memperoleh klasifikasi biner konjungtivitis yang presisi.

\section{Dasar Teori}
\label{sec:dasar-teori}

\subsection{Anatomi Mata}
\label{subsec:anatomi-mata}

Lapisan terluar bola mata manusia tersusun atas kornea di bagian anterior dan sklera di bagian posterior \parencite{wandini2024konjungtivitis}. Sklera berfungsi sebagai dinding pelindung bola mata yang berwarna putih opak dan tersusun atas anyaman padat jaringan fibrosa kolagen \parencite{wandini2024konjungtivitis}. Area sklera anterior dilapisi oleh membran mukosa tipis dan transparan yang disebut konjungtiva bulbi \parencite{sitompul2017konjungtivitis}. Konjungtiva ini longgar dan memiliki vaskularisasi mikrosirkulasi yang kaya, yang terletak tepat di atas warna dasar putih dari sklera \parencite{wandini2024konjungtivitis}. Kehadiran latar belakang putih dari jaringan sklera ini sangat menguntungkan secara klinis karena memberikan kontras warna yang tinggi untuk mendeteksi dilatasi pembuluh darah konjungtiva \parencite{sitompul2017konjungtivitis, wandini2024konjungtivitis}.

Pada penelitian kali ini, teori ini digunakan untuk memahami struktur anatomi sklera anterior dan konjungtiva bulbi sebagai area target pengamatan utama pada penelitian ini.

\subsection{Konjungtivitis}
\label{subsec:konjungtivitis}

Konjungtivitis didefinisikan sebagai peradangan atau inflamasi pada lapisan konjungtiva yang disebabkan oleh infeksi bakteri, virus, reaksi alergi, maupun paparan zat toksik \parencite{fhadila2023identifikasi, ramadirta2023konjungtivitis}. Secara klinis, proses peradangan ini memicu respons imun berupa dilatasi pembuluh darah (hiperemia), pembengkakan jaringan (kemosis), serta produksi eksudat \parencite{sitompul2017konjungtivitis, fhadila2023identifikasi}. Gejala hiperemia vaskular ini menyebabkan mata tampak merah karena aliran darah mikro meningkat pesat pada pembuluh darah konjungtiva bulbi \parencite{wandini2024konjungtivitis}. Deteksi dini tanda kemerahan ini sangat penting untuk mencegah penularan konjungtivitis menular dan membedakannya dari penyakit okular berbahaya lainnya \parencite{sitompul2017konjungtivitis, ramadirta2023konjungtivitis}.

Pada penelitian kali ini, teori ini digunakan untuk mengenali karakteristik visual patologis berupa warna kemerahan dan pelebaran pembuluh darah konjungtiva sebagai indikator utama penyakit konjungtivitis pada penelitian ini.

\subsection{Segmentasi Sklera}
\label{subsec:segmentasi-sklera}

Dalam analisis citra medis anterior mata secara otomatis, kehadiran objek non-okular di luar sklera (seperti kelopak mata, bulu mata, iris, pupil, dan kulit wajah) bertindak sebagai derau (\textit{noise}) yang dapat menurunkan akurasi model klasifikasi global \parencite{sangeethaa2024enhanced, jiang2023automatic}. Untuk meminimalkan bias tersebut, segmentasi sklera digunakan untuk mendeteksi dan mengisolasi wilayah sklera putih secara spesifik \parencite{madhusudhan2026deep}. Segmentasi semantik berbasis model pembelajaran mendalam (seperti arsitektur U-Net atau U-Net++) terbukti mampu melakukan penandaan piksel sklera dari gambar anterior mata \parencite{madhusudhan2026deep, fuentes2025illumination}. Hasil prediksi masker biner sklera tersebut kemudian diaplikasikan pada citra asli untuk meniadakan informasi di luar sklera (mengubahnya menjadi hitam atau bernilai nol), sehingga menyisakan citra bersih yang siap diklasifikasikan \parencite{fuentes2025illumination}.

Pada penelitian kali ini, teori ini digunakan untuk memotong dan mengisolasi objek sklera dari objek non-okular di sekitarnya agar model klasifikasi hanya berfokus menganalisis citra sklera bersih tanpa terganggu oleh derau eksternal pada penelitian ini.

\subsection{Fitur Vaskular}
\label{subsec:fitur-vaskular}

Fitur vaskular merepresentasikan geometri, percabangan, densitas, serta intensitas warna dari jalinan pembuluh darah yang melintasi konjungtiva bulbi di atas sklera anterior \parencite{madhusudhan2026deep}. Pembuluh darah ini sangat rentan terdistorsi oleh pantulan cahaya (\textit{glare}) permukaan basah mata dan variasi pencahayaan luar ruangan \parencite{fuentes2025illumination}. Prapemrosesan citra dilakukan dengan restorasi wilayah \textit{glare} pada ruang warna CIE LAB melalui algoritma \textit{inpainting} serta aplikasi CLAHE untuk mempertegas garis pembuluh darah lokal \parencite{fuentes2025illumination, madhusudhan2026deep}. Pengukuran keparahan dilatasi pembuluh darah ini dapat dirumuskan melalui indeks kemerahan pembuluh darah (\textit{Redness Index}) dan indeks kepadatan pembuluh darah (\textit{Vessel Density Index}). Secara umum, \textit{Redness Index} ($RI$) pada piksel ke-$i$ dirumuskan pada Persamaan \ref{eq:redness-index}:
\begin{equation}
  \label{eq:redness-index}
  RI_i = \frac{R_i}{R_i + G_i + B_i}
\end{equation}
di mana $R_i$, $G_i$, dan $B_i$ menyatakan nilai intensitas saluran warna merah, hijau, dan biru pada piksel tersebut \parencite{madhusudhan2026deep}. Sedangkan \textit{Vessel Density Index} ($VDI$) dihitung sebagai rasio luas piksel pembuluh darah merah hasil ekstraksi terhadap luas total wilayah sklera \parencite{fuentes2025illumination}.

Pada penelitian kali ini, teori ini digunakan untuk menormalisasi variasi kecerahan gambar, meminimalkan pengaruh pantulan cahaya, serta mengekstraksi dan memperjelas pola garis-garis pembuluh darah sklera yang mengalami inflamasi pada penelitian ini.

\subsection{CNN}
\label{subsec:cnn}

\textit{Convolutional Neural Network} (CNN) merupakan salah satu arsitektur \textit{deep learning} paling populer untuk pemrosesan citra digital yang secara otomatis mengekstrak fitur spasial hierarkis dari citra okular \parencite{ali2025automated}. Jaringan CNN yang dalam seperti ResNet menggunakan koneksi pemintas (\textit{skip connection} atau \textit{shortcut connection}) untuk melewatkan peta fitur langsung antar lapisan tanpa mengalami penurunan gradien (\textit{vanishing gradient}) selama pelatihan \parencite{sangeethaa2024enhanced}. Fitur-fitur visual lokal seperti warna merah, tekstur serat pembuluh darah, dan intensitas dilatasi diekstraksi melalui lapisan-lapisan konvolusi hierarkis untuk menghasilkan representasi fitur yang kokoh sebelum dilakukan klasifikasi kelas \parencite{ali2025automated, sangeethaa2024enhanced}.

Pada penelitian kali ini, teori ini digunakan sebagai landasan arsitektur dasar jaringan saraf konvolusional (\textit{backbone}) untuk mengekstraksi fitur visual dari citra sklera pada penelitian ini.

\subsection{Attention Mechanism}
\label{subsec:attention-mechanism}

Mekanisme perhatian (\textit{attention mechanism}) dalam arsitektur CNN dirancang untuk memberikan bobot relevansi dinamis pada peta fitur hasil konvolusi agar model dapat memusatkan pembelajaran pada area target yang paling informatif \parencite{jiang2023automatic, sangeethaa2024enhanced}. Mekanisme perhatian spasial (\textit{spatial attention}) berfokus pada letak spasial (\textit{where}) wilayah pembuluh darah yang melebar pada sklera \parencite{jiang2023automatic}. Secara matematis, \textit{Spatial Attention} ($\mathbf{M}_s$) dihitung berdasarkan operasi penggabungan (\textit{concatenation}) rata-rata fitur dan nilai maksimum fitur sepanjang dimensi saluran warna, dilanjutkan dengan konvolusi filter $7 \times 7$ dan fungsi aktivasi sigmoid pada Persamaan \ref{eq:spatial-attention}:
\begin{equation}
  \label{eq:spatial-attention}
  \mathbf{M}_s(\mathbf{F}) = \sigma\left(f^{7 \times 7}\left(\left[\text{AvgPool}(\mathbf{F}); \text{MaxPool}(\mathbf{F})\right]\right)\right)
\end{equation}
di mana $\mathbf{F}$ menyatakan peta fitur masukan dari lapisan konvolusi, $\text{AvgPool}(\mathbf{F})$ menyatakan rata-rata fitur sepanjang dimensi saluran, $\text{MaxPool}(\mathbf{F})$ menyatakan nilai maksimum fitur sepanjang dimensi saluran, $[\cdot ; \cdot]$ menyatakan operasi penggabungan (\textit{concatenation}), $f^{7 \times 7}$ menyatakan operasi konvolusi dengan ukuran filter $7 \times 7$, dan $\sigma$ menyatakan fungsi aktivasi sigmoid \parencite{sangeethaa2024enhanced, jiang2023automatic}. Hasil perkalian spasial ($\mathbf{M}_s(\mathbf{F}) \otimes \mathbf{F}$) ini meminimalkan bias fokus model akibat daerah sklera sehat yang berwarna putih bersih \parencite{sangeethaa2024enhanced}.

Pada penelitian kali ini, teori ini digunakan untuk mengintegrasikan modul perhatian spasial pada jaringan CNN agar model lebih fokus pada wilayah lesi pembuluh darah yang meradang pada penelitian ini.

\subsection{Imbalanced Dataset}
\label{subsec:imbalanced-dataset}

Ketidakseimbangan data (\textit{class imbalance}) terjadi ketika jumlah sampel pada kelas mayoritas (Normal) jauh lebih melimpah dibandingkan sampel pada kelas minoritas (Konjungtivitis), yang berisiko membuat klasifikasi model menjadi bias mengabaikan kelas minoritas \parencite{ali2025investigating, jiang2023automatic}. Dua teknik penanganan ketidakseimbangan kelas ini meliputi manipulasi data dan penyesuaian fungsi kerugian \parencite{ali2025investigating, fuentes2025illumination}. Augmentasi data geometris (seperti rotasi, pembalikan horizontal/vertikal, dan pergeseran piksel) dilakukan khusus untuk melipatgandakan sampel kelas konjungtivitis secara sintetis \parencite{ali2025automated}. Selain itu, \textit{Class Weighting} diterapkan pada fungsi kerugian \textit{Binary Cross-Entropy} untuk memberikan penalti kerugian yang lebih berat saat model salah mengklasifikasikan kelas minoritas \parencite{jiang2023automatic, ali2025investigating}.

Pada penelitian kali ini, teori ini digunakan untuk merancang strategi augmentasi dan fungsi kerugian berbobot guna melatih model secara adil pada dataset yang tidak seimbang pada penelitian ini.

\subsection{Metrik Evaluasi}
\label{subsec:metrik-evaluasi}

Pengukuran kinerja model klasifikasi bertujuan untuk menguji kemampuan generalisasi arsitektur usulan pada data uji independen \parencite{ali2025investigating}. Mengingat dataset medis yang digunakan memiliki karakteristik sangat timpang (\textit{class imbalance}), penggunaan metrik evaluasi tunggal seperti akurasi (\textit{accuracy}) dapat menyesatkan karena cenderung mencerminkan kelas mayoritas \parencite{ali2025investigating, jiang2023automatic}. Oleh karena itu, evaluasi model wajib menggunakan kombinasi metrik yang meliputi \textit{Confusion Matrix}, akurasi (\textit{accuracy}), presisi (\textit{precision}), sensitivitas (\textit{recall}), dan rata-rata harmonik keduanya (\textit{F1-Score}) \parencite{ali2025automated, fuentes2025illumination}.

Secara matematis, metrik-metrik evaluasi tersebut didefinisikan berdasarkan nilai \textit{True Positive} ($TP$), \textit{True Negative} ($TN$), \textit{False Positive} ($FP$), dan \textit{False Negative} ($FN$). Akurasi merepresentasikan persentase kecocokan prediksi model terhadap seluruh data aktual, yang dirumuskan pada Persamaan \ref{eq:akurasi}:
\begin{equation}
  \label{eq:akurasi}
  \text{Akurasi} = \frac{TP + TN}{TP + TN + FP + FN}
\end{equation}

Presisi menunjukkan kemampuan model dalam meminimalkan kesalahan prediksi positif palsu (\textit{false alarm}), yang dirumuskan pada Persamaan \ref{eq:presisi}:
\begin{equation}
  \label{eq:presisi}
  \text{Presisi} = \frac{TP}{TP + FP}
\end{equation}

Sensitivitas (\textit{recall}) menunjukkan kemampuan model dalam mengenali kembali seluruh kasus positif aktual tanpa terlewat, yang dirumuskan pada Persamaan \ref{eq:recall}:
\begin{equation}
  \label{eq:recall}
  \text{Recall} = \frac{TP}{TP + FN}
\end{equation}

Sedangkan \textit{F1-Score} merupakan rata-rata harmonik antara presisi dan \textit{recall} untuk menguji performa model secara seimbang pada data timpang, yang dirumuskan pada Persamaan \ref{eq:f1score}:
\begin{equation}
  \label{eq:f1score}
  \text{F1-Score} = 2 \times \frac{\text{Presisi} \times \text{Recall}}{\text{Presisi} + \text{Recall}}
\end{equation}

Pada penelitian kali ini, teori ini digunakan untuk menilai dan memvalidasi keandalan serta akurasi klasifikasi model dalam mendeteksi konjungtivitis pada penelitian ini.
